\section{切—投影模型集与有限分辨率读出}\label{sec:cut-project-readout}

\subsection{模型集（D3.1）}

取紧致可测窗口（接受域）
\[
W\subset E_{\perp},
\]
并取内部相位（窗口平移）$u\in E_{\perp}$。定义切—投影点集
\[
\Lambda(W,u):=\left\{\pi_{\parallel}(n):n\in\mathcal{L},\ \pi_{\perp}(n)\in W+u\right\}\subset E_{\parallel}.
\]
若考虑标记，则对 $x=\pi_{\parallel}(n)\in\Lambda(W,u)$ 定义诱导标记
\[
\Sigma_{u}(x):=\sigma(n),
\]
其良定性由（H2.3）保证。

\subsection{核化窗口与复幅度读出（D3.2）}

理想窗口指示函数 $\ind_{W+u}$ 对应尖锐选择，难以表征有限分辨率。取核
\[
K_{\varepsilon}:E_{\perp}\to[0,\infty),\qquad \varepsilon>0,
\]
满足 $\int_{E_{\perp}}K_{\varepsilon}(y)\,\dd y=1$，且当 $\varepsilon\to 0$ 时 $K_{\varepsilon}$ 逼近 Dirac 族。取幅度映射
\[
a:\mathcal{A}\to\CC.
\]
定义在 $E_{\parallel}$ 上的复测度（或分布）
\[
\Psi_{\varepsilon,u}:=\sum_{n\in\mathcal{L}} a(\sigma(n))\,K_{\varepsilon}(\pi_{\perp}(n)-u)\,\delta_{\pi_{\parallel}(n)} .
\]
为避免分布平方的技术问题，可再引入物理空间点扩散核 $\phi_{\delta}$（$\delta>0$）并定义平滑场
\[
\psi_{\varepsilon,u}(x):=(\Psi_{\varepsilon,u}*\phi_{\delta})(x),
\]
相应强度读出定义为
\[
I_{\varepsilon,u}(x):=\big|\psi_{\varepsilon,u}(x)\big|^{2}.
\]
本文将“观测统计”理解为由 $I_{\varepsilon,u}$ 或其采样机制诱导的可重复统计量（例如相关函数、功率谱、事件计数等）。

\subsection{相干叠加项的结构来源（T3.3）}

对任意固定 $x\in E_{\parallel}$，若 $\psi_{\varepsilon,u}(x)$ 可写为有限或可控求和 $\sum_{j}A_j(x)$，则
\[
I_{\varepsilon,u}(x)=\left|\sum_{j}A_j(x)\right|^{2}
=\sum_{j}|A_j(x)|^{2}+\sum_{j\neq k}A_j(x)\overline{A_k(x)}.
\]
因此，干涉项的存在由两类机制决定：
\begin{itemize}
  \item 多个投影贡献在物理空间核 $\phi_{\delta}$ 的支撑内叠加；
  \item 多个内部坐标在 $K_{\varepsilon}$ 的有效支撑内同时获得非零权重。
\end{itemize}
该结论属于代数恒等式层面（T），不依赖任何语义映射。

